\section{The full chromatic number in dimension six}
\label{sec:chromatic}

\begin{theorem}\label{thm:chromatic}
\(\chi(\OG_6^*)=15\).
\end{theorem}

\begin{proof}
Let \(T=\Span(3,12,48)\) in binary integer coordinates, with bit
\(i-1\) representing the coefficient of \(e_i\).  Thus its generators
are \(e_1+e_2,\ e_3+e_4,\ e_5+e_6\).  The form vanishes on \(T\), so
its seven lines, seven planes and \(T\) itself form a \(15\)-clique.
Hence \(\chi(\OG_6^*)\ge15\).

For the upper bound, the file
\path{results/full-15-coloring.json} gives a basis and a color in
\(\{0,\ldots,14\}\) for every nonzero subspace.  The independent
checker reconstructs each span and verifies basis independence,
distinctness, and the dimension counts
\[
63,\ 651,\ 1395,\ 651,\ 63,\ 1 .
\]
These equal the Gaussian binomial coefficients for \(d=1,\dots,6\),
which proves completeness of the catalogue.  For all \(3\,986\,076\)
pairs of distinct entries the checker tests orthogonality on the
basis vectors.  Bilinearity identifies this test with graph
adjacency.  All \(44\,968\) orthogonality edges have distinct
endpoint colors.  This establishes the upper bound
\(\chi(\OG_6^*)\le15\), and therefore equality.
\end{proof}

\begin{remark}\label{rem:precoloring}
To find the certificate, fix distinct colors on the \(15\) nonzero
subspaces \(W\) of \(T\).  Every \(15\)-coloring uses all colors on
this clique, so relabeling colors makes this precoloring an
equivalent instance.  The forbidden colors for a vertex
\(U\not\subseteq T\) are exactly those \(W\subseteq H(U)=U^\perp\cap
T\).  If \(H(U)=T\), then \(U\subseteq T^\perp=T\), a contradiction.
Thus \(\dim H(U)\in\{0,1,2\}\), giving lists of sizes
\[
15-N(\dim H(U))\in\{15,14,11\} .
\]
\end{remark}

\begin{remark}\label{rem:degree-deletion}
A vertex whose residual degree is smaller than its list size can be
deleted and restored greedily after the remaining graph is colored.
Repeated deletion removes \(729\) uncolored vertices and leaves
\(2\,080\) vertices with \(39\,400\) edges.  The list-coloring
encoding has \(28\,600\) variables and \(627\,510\) clauses.  There
is one Boolean variable for each allowed vertex--color pair.  Each
vertex has an at-least-one clause and pairwise at-most-one clauses;
each edge forbids simultaneous use of any color common to its
endpoints' lists.  Hence satisfying assignments are exactly the
proper list colorings of the residual graph.  A satisfying
assignment was extended in reverse deletion order to produce the
full certificate checked above.
\end{remark}
